\documentclass[11pt]{article}
\usepackage[T1]{fontenc}
\usepackage[utf8]{inputenc}
\usepackage{lmodern}
\usepackage[margin=1in]{geometry}
\usepackage{amsmath,amssymb,amsthm,mathtools}
\usepackage{enumitem}
\usepackage{hyperref}

\hypersetup{colorlinks=true,linkcolor=blue,citecolor=blue,urlcolor=blue}

\newtheorem{theorem}{Theorem}[section]
\newtheorem{corollary}[theorem]{Corollary}
\newtheorem{definition}[theorem]{Definition}
\newtheorem{assumption}[theorem]{Assumption}

\title{Reality as a Consensus Protocol:\\
The Fixed-Point Computation That Implements Physics}
\author{B.\ M\"uller}
\date{March 2026}

\begin{document}
\maketitle

\begin{abstract}
We show that objective physical law is the unique fixed point of a distributed consensus protocol running between observer patches. The model is simple: a finite graph of local observers, each holding a local state, connected by overlap regions where they must agree. No global state is postulated. Instead, an asynchronous reconciliation process repairs local inconsistencies, and we prove that if this process terminates and is locally confluent, every initial configuration converges to a unique normal form regardless of execution order. The terminal state \emph{is} physical reality.

We then prove that pairwise agreement between neighboring observers is not enough to guarantee a globally consistent reality. A cycle-holonomy obstruction shows that topology can permanently prevent global reconciliation. The minimal example is a triangle over $\mathbb{Z}_2$ where every edge looks fine but the whole system is frustrated. These irreducible frustrations are stable defects of the reconciliation process, and we identify them with particles.

The framework yields four exact theorems: (1) asynchronous confluence, giving schedule-independent physical law; (2) a cycle-obstruction criterion showing when local consistency fails globally; (3) a gauge quotient theorem identifying gauge symmetry with implementation hiding; and (4) a record-layer convergence theorem giving CRDT-like semantics for classical observations. We also formalize a selection principle over the space of possible reconciliation laws, where schedule-robust, observer-supporting, and simple laws are favored by a monotonically increasing fitness functional.

The formalism is the computational spine of Observer-Patch Holography (OPH), a framework in which a conditional scaling-limit Einstein branch, compact gauge reconstruction, an extended-MAR Standard Model closure, and the string-theoretic worldsheet expansion are developed as emergent sectors. The full OPH manuscript is available at \url{https://oph-book.floatingpragma.io/papers/observers_are_all_you_need.pdf}.
\end{abstract}

\section{Introduction}

Here is the claim, stated as plainly as possible:

\begin{center}
\fbox{\parbox{0.88\textwidth}{\centering
Objective physical law $=$ the gauge-invariant normal form\\
of an asynchronous local reconciliation process.}}
\end{center}

That becomes mathematically sharp once you define patches, overlaps, local repair maps, and gauge equivalence. By the end of this paper you will have the proofs.

The idea goes like this. Suppose the universe has no global state. No God's-eye view, no master clock, no centralized ledger. All that exists are local observers, each holding a local description of their patch of reality, connected to neighbors through overlap regions where their descriptions must be compatible. When two neighbors disagree on their shared boundary, a local repair rule fires and updates one of them. This process runs asynchronously, with no preferred order of execution.

The question is: does this process converge? And if it converges, does the answer depend on the order in which repairs happen?

If the answer to both questions is yes, then the terminal state of the reconciliation process \emph{is} an objective physical law. It does not matter who ran the updates first, or which observer got corrected before which other observer. The outcome is the same. That is a mathematical theorem, not a philosophical position.

Anyone who has built distributed systems will recognize the structure immediately. This is a consistency protocol. The local repair maps are like conflict-resolution rules in a CRDT. The overlap projections are like shared keys in a distributed database. The convergence theorem is an eventual-consistency guarantee. And gauge symmetry is just the freedom to choose different internal representations on each node, as long as the interface data stays the same. In a distributed system, that is called encapsulation. In physics, it is called local gauge invariance. They are the same thing.

This paper proves four core results:

\begin{enumerate}
\item \textbf{Asynchronous confluence.} Under termination and local confluence, every initial state has a unique normal form, independent of update schedule (Theorem~\ref{thm:confluence}).
\item \textbf{Cycle obstruction.} For affine overlap constraints over an abelian group, global consistency holds if and only if the holonomy vanishes on every cycle (Theorem~\ref{thm:holonomy}). The parity triangle gives the minimal frustrated example.
\item \textbf{Gauge quotient.} Under gauge-covariant dynamics, the normal-form map descends to the quotient. Gauge-invariant observables are strictly unique (Theorem~\ref{thm:gauge}).
\item \textbf{Record-layer convergence.} Commuting closure operators on a finite lattice converge to the least common fixed point under any fair asynchronous schedule (Theorem~\ref{thm:crdt}).
\end{enumerate}

We also define a fitness functional over the space of candidate reconciliation laws and prove that replicator dynamics monotonically increases mean fitness (Theorem~\ref{thm:replicator}). This gives a clean mathematical model for the selection of physical law.

The results here are exact theorems about a computational model. The further claim that physical reality instantiates this model is the interpretive hypothesis of Observer-Patch Holography (OPH). The full OPH framework, from which general relativity, the SM gauge group, and the string-theory worldsheet expansion are derived, is presented in a companion manuscript~\cite{OPH}.


\section{Patch Nets, Overlaps, and Global Consistency}

\begin{definition}[Patch net]\label{def:patchnet}
Let $G=(V,E)$ be a finite connected graph. Each vertex $i\in V$ is an \textbf{observer patch} with finite local state space $S_i$. The global state space is
\[
\Sigma := \prod_{i\in V} S_i.
\]
For each edge $e=\{i,j\}\in E$, let $I_e$ be an interface alphabet and let
\[
\pi_{i,e}:S_i\to I_e,
\qquad
\pi_{j,e}:S_j\to I_e
\]
be the interface projection maps. A global state $s=(s_i)_{i\in V}\in\Sigma$ is \textbf{consistent on edge $e=\{i,j\}$} iff
\[
\pi_{i,e}(s_i)=\pi_{j,e}(s_j).
\]
The global consistency set is
\[
C:=\bigl\{s\in\Sigma:\forall\, e=\{i,j\}\in E,\ \pi_{i,e}(s_i)=\pi_{j,e}(s_j)\bigr\}.
\]
\end{definition}

For exposition we use a finite pairwise-overlap graph. The hypergraph version is straightforward: replace edges by hyperedges and pairwise equality by a common interface label on each hyperedge. Nothing in the proofs depends on the pairwise restriction.

The picture: each observer holds a local state, and neighboring observers share an interface through which they can compare notes. A universe-state is physically admissible exactly when all neighbors agree on their shared data. This is a constraint satisfaction problem (CSP), and the consistent states are the codewords.

\begin{definition}[Inconsistency potential]\label{def:potential}
For each edge $e$, choose a weight $w_e>0$ and a function $d_e:I_e\times I_e\to\mathbb{R}_{\ge 0}$ with $d_e(a,b)=0 \iff a=b$. Define
\[
\Phi(s)
:=
\sum_{e=\{i,j\}\in E}
w_e\, d_e\!\bigl(\pi_{i,e}(s_i),\pi_{j,e}(s_j)\bigr).
\]
Then $s\in C \iff \Phi(s)=0$.
\end{definition}

So $\Phi$ is the total disagreement energy of the universe. Consistent states have zero energy. Everything else is frustrated.


\section{Asynchronous Reconciliation and the Main Fixed-Point Theorem}

\begin{definition}[Local repair law]\label{def:repair}
A \textbf{law} $\lambda$ is a family of local repair maps
\[
T_i^\lambda:\Sigma\to\Sigma
\qquad (i\in V)
\]
such that $T_i^\lambda$ changes only the state of patch $i$ (or, more generally, only a bounded neighborhood of $i$). Write $s\to_i t$ iff $t=T_i^\lambda(s)\neq s$. Let $\to := \bigcup_{i\in V}\to_i$, and let $\to^*$ be its reflexive-transitive closure. A state $s\in\Sigma$ is a \textbf{normal form} iff $T_i^\lambda(s)=s$ for all $i\in V$.
\end{definition}

The three assumptions that make the machine work:

\begin{assumption}[Repair completeness]\label{ass:complete}
$s\in C \iff \forall\, i\in V,\; T_i^\lambda(s)=s.$
\end{assumption}

Normal forms are exactly the globally consistent states. The dynamics is neither too weak (missing some inconsistencies) nor too strong (repairing things that were fine).

\begin{assumption}[Termination]\label{ass:termination}
Every enabled repair strictly decreases the inconsistency potential: $s\to t \implies \Phi(t)<\Phi(s)$.
\end{assumption}

Since $\Sigma$ is finite, this guarantees termination.

\begin{assumption}[Local confluence]\label{ass:confluence}
The repair relation is locally confluent: if $s\to t$ and $s\to u$, then there exists $v\in\Sigma$ with $t\to^* v$ and $u\to^* v$.
\end{assumption}

This is the schedule-independence assumption written as a diamond property.

\begin{theorem}[Asynchronous confluence / fixed-point law]\label{thm:confluence}
Under Assumptions~\ref{ass:complete}--\ref{ass:confluence}, every initial state $s\in\Sigma$ has a unique normal form
\[
\operatorname{nf}_\lambda(s)\in C,
\]
and every maximal asynchronous repair execution from $s$ terminates at that same state. The terminal state is independent of update order.
\end{theorem}

\begin{proof}
By Assumption~\ref{ass:termination}, every rewrite step strictly decreases $\Phi$, so no infinite rewrite sequence exists. Thus $\to$ is terminating.

We prove confluence by induction on $\Phi(s)$. Take any $s\in\Sigma$ and suppose $s\to^* t$ and $s\to^* u$. We must show $t$ and $u$ have a common descendant.

If one reduction has length~0, there is nothing to prove. Otherwise write $s\to s_1\to^* t$ and $s\to s_2\to^* u$. By local confluence, there exists $w$ with $s_1\to^* w$ and $s_2\to^* w$. Since $s\to s_1$ and $s\to s_2$, Assumption~\ref{ass:termination} gives $\Phi(s_1)<\Phi(s)$ and $\Phi(s_2)<\Phi(s)$. The induction hypothesis applies at $s_1$ and $s_2$. Hence $t$ and $w$ have a common descendant, and $u$ and $w$ have a common descendant. Therefore $t$ and $u$ have a common descendant.

A terminating confluent rewrite system has a unique normal form reachable from each initial state. By Assumption~\ref{ass:complete}, the normal forms are exactly $C$. Every maximal execution terminates (by $\Phi$-decrease) and reaches $\operatorname{nf}_\lambda(s)$ (by uniqueness).
\end{proof}

\begin{corollary}[Objective law is schedule-independent]\label{cor:objective}
Let $M:\Sigma\to Y$ be any observable. Under the hypotheses of Theorem~\ref{thm:confluence}, $M(\operatorname{nf}_\lambda(s))$ is independent of the asynchronous update schedule. If physical law is identified with the map $s\mapsto M(\operatorname{nf}_\lambda(s))$, then physical law is objective.
\end{corollary}

\begin{proof}
All schedules from the same initial $s$ terminate at $\operatorname{nf}_\lambda(s)$, so all yield the same $M$-value.
\end{proof}

This is the main theorem of the paper, and it deserves a box:
\[
\boxed{\text{termination + local confluence} \;\Longrightarrow\; \text{unique schedule-independent normal form}}
\]

The universe does not need a master clock or a preferred reference frame. As long as the reconciliation process terminates and is locally confluent, every observer, running repairs in whatever order they happen to run them, arrives at the same answer. Objectivity is a convergence guarantee.


\section{Why Local Agreement Is Not Enough: Cycle Holonomy and Frustration}

Now we show that the story has a twist: just because every pair of neighbors agrees does not mean the whole system is consistent.

\begin{theorem}[Cycle-obstruction / holonomy criterion]\label{thm:holonomy}
Let $A$ be an abelian group, and let $G=(V,E)$ be a connected graph with an arbitrary orientation on each edge. For each oriented edge $e:u\to v$, assign a label $b_e\in A$. Consider the affine consistency equations
\[
x_v - x_u = b_e
\qquad
\text{for every oriented edge } e:u\to v,
\]
where $x_v\in A$ are unknown patch labels. A global solution $x:V\to A$ exists if and only if for every cycle $C\subseteq G$,
\[
\sum_{e\in C}\varepsilon_C(e)\, b_e = 0,
\]
where $\varepsilon_C(e)=+1$ if the cycle traverses $e$ in the chosen orientation and $-1$ otherwise.
\end{theorem}

\begin{proof}
\textbf{Necessity.} Suppose $x$ is a solution. Summing the edge equations around any cycle $C$,
\[
\sum_{e:u\to v\in C}\varepsilon_C(e)\,(x_v-x_u)
=
\sum_{e\in C}\varepsilon_C(e)\,b_e.
\]
The left side telescopes to $0$ because every vertex appears once with $+$ sign and once with $-$ sign.

\textbf{Sufficiency.} Fix a root $r\in V$. For any vertex $v$, choose a path $P_{r\to v}$ and define
\[
x_v := \sum_{e\in P_{r\to v}} \varepsilon_{P_{r\to v}}(e)\, b_e,
\qquad x_r:=0.
\]
If $P$ and $P'$ are two paths from $r$ to $v$, traversing $P$ followed by the reverse of $P'$ yields a cycle. By the vanishing-holonomy assumption, the total signed sum is zero, so $x_v$ is well-defined. For any edge $e:u\to v$, extending a path to $u$ by that edge gives $x_v = x_u + b_e$.
\end{proof}

\begin{corollary}[Parity triangle: pairwise consistency is not enough]\label{cor:triangle}
Take $A=\mathbb{Z}_2$ on the triangle $A$-$B$-$C$-$A$ with edge labels $b_{AB}=0$, $b_{BC}=0$, $b_{CA}=1$. Each individual edge equation is satisfiable. But the global system is not: the cycle sum is $0\oplus 0\oplus 1 = 1\neq 0$.
\end{corollary}

This is the simplest possible demonstration that local agreement does not imply global reality. Every pair of neighbors can find states that work for them, but there is no assignment that works for everyone simultaneously. The obstruction lives on the cycle.

\begin{corollary}[Stable defects as frustrated holonomy]\label{cor:defects}
Define the defect energy
\[
\Phi_b(x) := \sum_{e:u\to v\in E} w_e\,\mathbf{1}\!\left[x_v - x_u \neq b_e\right].
\]
If the cycle-holonomy condition fails, then $\min_{x:V\to A}\Phi_b(x)>0$. Every minimizer contains irreducible residual inconsistency.
\end{corollary}

\begin{proof}
If $\min_x\Phi_b(x)=0$, some assignment satisfies all edge equations, contradicting Theorem~\ref{thm:holonomy}.
\end{proof}

And here is the physical punchline: these residual inconsistencies cannot be repaired. They are topologically protected. The reconciliation process pushes them around but can never eliminate them. They are stable, localized concentrations of frustrated holonomy.

In this framework, that is what a particle is. A particle is a defect in the reconciliation process that cannot be smoothed away because the topology forbids it.


\section{Gauge Symmetry as Implementation Hiding}

\begin{definition}[Gauge action]\label{def:gauge}
Let $\Gamma = \prod_{i\in V}\Gamma_i$ act on $\Sigma$ componentwise. The action is a \textbf{gauge action} if for every $e=\{i,j\}\in E$,
\[
\pi_{i,e}(\gamma_i\cdot x)=\pi_{i,e}(x)
\qquad
\forall\, x\in S_i,\ \forall\, \gamma_i\in \Gamma_i.
\]
Gauge changes alter hidden local representations but do not alter overlap data. Assume gauge covariance:
\[
T_i^\lambda(\gamma\cdot s)=\gamma\cdot T_i^\lambda(s)
\qquad
\forall\, i,\ \forall\, \gamma\in\Gamma,\ \forall\, s\in\Sigma.
\]
\end{definition}

In software terms: each node can refactor its internal state however it likes, as long as the API it exposes to its neighbors stays the same. That is encapsulation. That is also gauge invariance.

\begin{theorem}[Gauge quotient theorem]\label{thm:gauge}
Under the gauge action of Definition~\ref{def:gauge} and the hypotheses of Theorem~\ref{thm:confluence},
\[
\operatorname{nf}_\lambda(\gamma\cdot s)
=
\gamma\cdot \operatorname{nf}_\lambda(s)
\qquad
\forall\, \gamma\in\Gamma,\ \forall\, s\in\Sigma.
\]
Hence the normal-form map descends to the quotient:
\[
\overline{\operatorname{nf}}_\lambda:\Sigma/\Gamma \to \Sigma/\Gamma,
\qquad
[s]\mapsto [\operatorname{nf}_\lambda(s)].
\]
\end{theorem}

\begin{proof}
Suppose $s\to_i t$, so $t=T_i^\lambda(s)\neq s$. Gauge covariance gives $T_i^\lambda(\gamma\cdot s) = \gamma\cdot T_i^\lambda(s) = \gamma\cdot t$, so $\gamma\cdot s\to_i \gamma\cdot t$. The gauge action sends rewrite sequences to rewrite sequences: $s\to^* t$ implies $\gamma\cdot s\to^* \gamma\cdot t$.

Let $t=\operatorname{nf}_\lambda(s)$. Then $T_i^\lambda(t)=t$ for all $i$, so $T_i^\lambda(\gamma\cdot t) = \gamma\cdot T_i^\lambda(t) = \gamma\cdot t$, meaning $\gamma\cdot t$ is also normal. Since $\gamma\cdot s$ rewrites to $\gamma\cdot t$, uniqueness gives $\operatorname{nf}_\lambda(\gamma\cdot s) = \gamma\cdot\operatorname{nf}_\lambda(s)$.
\end{proof}

\begin{corollary}[Gauge-invariant law]\label{cor:gaugeinv}
If $M:\Sigma\to Y$ is gauge-invariant ($M(\gamma\cdot s)=M(s)$ for all $\gamma$), then $M(\operatorname{nf}_\lambda(s))$ depends only on the gauge orbit $[s]$, not the representative.
\end{corollary}

\begin{proof}
$M(\operatorname{nf}_\lambda(\gamma\cdot s)) = M(\gamma\cdot\operatorname{nf}_\lambda(s)) = M(\operatorname{nf}_\lambda(s))$.
\end{proof}

Uniqueness is physically meaningful only up to gauge. The universe's internal representation is not unique, but every measurement you can actually perform gives a unique answer. This is exactly what physicists mean when they say gauge symmetry is ``not a real symmetry.'' In our language: it is the freedom to choose different internal data structures on each node without affecting the distributed computation's output.


\section{Classical Records as an Eventually Consistent Layer}

\begin{definition}[Record lattice]\label{def:lattice}
Let $L$ be a finite distributive lattice. A map $F:L\to L$ is a \textbf{closure operator} if it is inflationary ($x\le F(x)$), monotone ($x\le y\implies F(x)\le F(y)$), and idempotent ($F(F(x))=F(x)$). Let $F_1,\dots,F_n$ be pairwise commuting closure operators on $L$.
\end{definition}

\begin{theorem}[Asynchronous convergence of the record layer]\label{thm:crdt}
For any initial record state $x_0\in L$, define $x^\star := F_1 F_2\cdots F_n(x_0)$. Then:
\begin{enumerate}
\item $x^\star$ is independent of the order of composition.
\item $x^\star$ is a common fixed point of all $F_i$.
\item $x^\star$ is the least common fixed point above $x_0$.
\item Any fair asynchronous application of the $F_i$ converges in finitely many steps to $x^\star$.
\end{enumerate}
\end{theorem}

\begin{proof}
Commutativity makes the composition order-independent, so $x^\star$ is well-defined.

For any $i$, commutativity and idempotence give $F_i(x^\star) = F_i F_1\cdots F_n(x_0) = F_1\cdots F_i^2\cdots F_n(x_0) = x^\star$. So $x^\star$ is a common fixed point.

If $y\ge x_0$ satisfies $F_i(y)=y$ for all $i$, monotonicity gives $x^\star = F_1\cdots F_n(x_0) \le F_1\cdots F_n(y) = y$. So $x^\star$ is the least common fixed point above $x_0$.

Under any fair schedule, once each operator has been applied at least once, the total effect is $F_1\cdots F_n$ regardless of order and repetitions (by commutativity and idempotence). The state reaches $x^\star$ and stays there.
\end{proof}

This is the record layer of reality. Classical observations, the stuff you can write down and share with other observers, form a lattice of shareable records. The reconciliation operators are closure operators. The classical world is the order-independent closure of the shareable record layer.

If you know distributed systems, this is a state-based CRDT with a join-semilattice. The theorem says the merge function is commutative, associative, and idempotent. Eventual consistency is guaranteed.


\section{Law-Space Selection and Observer Emergence}

Which reconciliation laws get selected? Not all laws are equally good. Some produce schedule-dependent outcomes (they fail the confluence condition). Some are schedule-robust but produce universes with no structure. We want laws that are robust, that support observers, and that are not needlessly complicated.

An important framing note: this section is a theorem about a \emph{meta-selection model}. We are not claiming that the actual universe literally ran replicator dynamics over law space. What we are doing is giving a clean mathematical statement of the selection idea and proving that it has the right monotonicity property. The interpretive step from here to cosmology is the paper's hypothesis, not its theorem.

We start by defining what it means for a law to support observers. An observer, operationally, is a \textbf{persistent predictive module}: a subgraph that maintains stable records and uses them to predict its boundary's future behavior.

\begin{definition}[Schedule robustness]\label{def:robustness}
Fix distributions $\mu$ over initial conditions and $\nu$ over asynchronous schedules, and a gauge-invariant observable $M$. For a law $\lambda$, define
\[
\mathcal{R}_M(\lambda)
:=
\Pr_{s\sim\mu,\;\sigma,\tau\sim\nu}
\!\left[
M\bigl(\operatorname{nf}^{\sigma}_\lambda(s)\bigr)
=
M\bigl(\operatorname{nf}^{\tau}_\lambda(s)\bigr)
\right].
\]
If Theorem~\ref{thm:confluence} holds for $\lambda$, then $\mathcal{R}_M(\lambda)=1$.
\end{definition}

\begin{definition}[Observer yield]\label{def:yield}
Let $X_t^\lambda$ denote the stationary process obtained by repeated local perturbation plus reconciliation under law $\lambda$. A subgraph $U\subseteq V$ with record map $r_U$ is $(\eta,\varepsilon,h)$-observer-like if it is record-stable:
\[
\Pr\bigl[r_U(X_{t+h}^\lambda)=r_U(X_t^\lambda)\bigr]\ge 1-\eta,
\]
and predictive:
\[
I\bigl(r_U(X_t^\lambda);\; X_{\partial U,\,t+1:t+h}^\lambda\bigr)\ge\varepsilon.
\]
Define
\[
\mathcal{O}_{\eta,\varepsilon,h}(\lambda)
:=
\mathbb{E}\!\left[
\#\left\{
U\subseteq V:
U \text{ is } (\eta,\varepsilon,h)\text{-observer-like}
\right\}
\right].
\]
\end{definition}

\begin{definition}[Law fitness]\label{def:fitness}
Let $K(\lambda)$ be a description-length penalty. Define
\[
f(\lambda) = \alpha\,\mathcal{R}_M(\lambda) + \beta\,\mathcal{O}_{\eta,\varepsilon,h}(\lambda) - \gamma\,K(\lambda),
\]
with $\alpha,\beta,\gamma>0$.
\end{definition}

\begin{theorem}[Replicator monotonicity on law space]\label{thm:replicator}
Let $\Lambda=\{\lambda_1,\dots,\lambda_m\}$ be candidate laws with population weights $x_i(t)$ under replicator dynamics:
\[
\dot{x}_i = x_i(f_i - \bar{f}),
\qquad
f_i := f(\lambda_i),
\qquad
\bar{f} := \sum_{j=1}^m x_j f_j.
\]
Then
\[
\frac{d}{dt}\bar{f} = \operatorname{Var}_x(f) \ge 0.
\]
Mean fitness is nondecreasing, and strictly increasing unless all extant laws have the same fitness.
\end{theorem}

\begin{proof}
Direct computation:
\[
\frac{d}{dt}\bar{f}
=
\sum_i \dot{x}_i f_i
=
\sum_i x_i(f_i - \bar{f})f_i
=
\sum_i x_i f_i^2 - \bar{f}^2
=
\operatorname{Var}_x(f) \ge 0.
\]
Equality iff all $f_i$ on the support of $x$ are equal.
\end{proof}

This is a theorem about the meta-model: a clean mathematical statement about how laws compete. Schedule-robust laws beat fragile ones. Observer-supporting laws beat barren ones. Simple laws beat overcomplicated ones. The variance of fitness drives selection. When all surviving laws have equal fitness, the process stops. That is equilibrium.


\section{Connection to Observer-Patch Holography}

The formalism above is the computational skeleton of Observer-Patch Holography (OPH). In OPH, the observer patches carry von Neumann algebras on a holographic screen $S^2$, the overlap projections are restrictions to shared subalgebras, and the consistency condition is algebraic state agreement on overlaps.

The bridge to physics works as follows. In the full OPH framework:

\begin{itemize}
\item The patch net becomes a net of subregion algebras on $S^2$.
\item Overlap consistency (Axiom A2 of OPH) is the algebraic version of Definition~\ref{def:patchnet}.
\item Local Markov recoverability (Axiom A4) provides the collar factorization $\rho_{ABD} = \bigoplus_\alpha p_\alpha\,\rho^{(\alpha)}_{Ab_L^\alpha}\otimes\rho^{(\alpha)}_{b_R^\alpha D}$ that makes observer checkpointing and restoration well-defined (see Appendix~\ref{app:splice}).
\item Gauge symmetry as implementation hiding (Theorem~\ref{thm:gauge}) becomes the edge-sector fusion structure from which OPH derives the Standard Model gauge group $\mathrm{SU}(3)\times\mathrm{SU}(2)\times\mathrm{U}(1)/\mathbb{Z}_6$.
\item Stable defects (Corollary~\ref{cor:defects}) become the topologically protected excitations that OPH identifies with particles.
\item The record-layer theorem (Theorem~\ref{thm:crdt}) provides the formal basis for the classical observation layer in OPH, where records are implemented by approximately commuting projectors in overlap centers.
\end{itemize}

The OPH manuscript derives general relativity from entanglement equilibrium and modular geometry, the complete SM gauge group from edge-sector admissibility, and the string-theoretic worldsheet expansion from large-$N$ heat-kernel asymptotics. These derivations are quantitative: they produce particle masses, coupling constants, and cosmological parameters from two input constants (the pixel area $a_{\rm cell}\approx 1.63\,\ell_P^2$ and the screen capacity $\log\dim\mathcal{H}_{\rm tot}\sim 10^{122}$).

The present paper provides the computational foundation. OPH says what the patches are, what the algebras are, and what the reconciliation dynamics looks like at the physical level. This paper proves that any system with the right abstract properties must converge to a unique, schedule-independent, gauge-invariant fixed point. That fixed point is what we call physical law.


\section{Discussion and Open Problems}

Five directions stand out as immediate next steps:

\begin{enumerate}
\item \textbf{Computational complexity.} What is the complexity of computing $\operatorname{nf}_\lambda(s)$? The termination bound from $\Phi$-decrease is a worst-case measure, but typical instances may converge much faster. Is there a polynomial bound for natural classes of patch nets?

\item \textbf{Coarse-graining.} When does reconciliation commute with renormalization? If you coarse-grain the patch net and then reconcile, do you get the same result as reconciling first and then coarse-graining? The conditions under which these operations commute would connect the discrete model to continuum field theory.

\item \textbf{Defect classification.} Which holonomy classes correspond to stable quasi-particles? In the abelian case, the classification is by the first cohomology group $H^1(G,A)$. The non-abelian case is richer and connects to the classification of topological defects in gauge theory.

\item \textbf{Observable-level confluence.} When the raw state is not unique (because local confluence fails), can gauge-invariant macrostates still be unique? This weaker form of objectivity may hold under broader conditions.

\item \textbf{Quantum lift.} How much of the patch-net theorem survives when local state spaces become operator algebras and repair maps become quantum channels? The Markov-collar splice theorem in the OPH appendix (where interior observer states can be spliced into different compatible environments without changing interior observables) suggests that the algebraic version is natural.
\end{enumerate}


\appendix
\section{Quantum/Algebraic Lift: Markov-Collar Splice Theorem}\label{app:splice}

This appendix ties the discrete computational model back to the algebraic OPH formalism. It is not required for the main CS theorems, but it shows that the patch-net structure is not floating free of the deeper framework.

In OPH, an observer is written as
\[
O=(P,\mathcal{A}(P),\rho,R),
\]
where $P$ is the screen patch, $\mathcal{A}(P)$ the local von Neumann algebra, $\rho$ the local state, and $R$ the record algebra.

\begin{theorem}[Markov-collar splice theorem]\label{thm:splice}
Suppose a collar tripartition $A$-$B$-$D$ has exact Markov decomposition
\[
\rho_{ABD}
=
\bigoplus_{\alpha}
p_\alpha\,
\rho^{(\alpha)}_{A b_L^\alpha}
\otimes
\rho^{(\alpha)}_{b_R^\alpha D}.
\]
Let $\sigma_{b_R^\alpha D'}^{(\alpha)}$ be any family of normalized environment states compatible with the same right-boundary sectors. Define
\[
\rho'_{AB D'}
=
\bigoplus_{\alpha}
p_\alpha\,
\rho^{(\alpha)}_{A b_L^\alpha}
\otimes
\sigma_{b_R^\alpha D'}^{(\alpha)}.
\]
Then for every observable $X$ supported on $A\cup b_L$,
\[
\operatorname{Tr}(X\rho'_{AB D'})
=
\operatorname{Tr}(X\rho_{ABD}).
\]
The interior observer state can be spliced into a different compatible environment without changing any interior observable statistics.
\end{theorem}

\begin{proof}
By blockwise factorization,
\[
\operatorname{Tr}(X\rho'_{AB D'})
=
\sum_\alpha
p_\alpha\,
\operatorname{Tr}\!\left(
X\, \rho^{(\alpha)}_{A b_L^\alpha}
\right)
\operatorname{Tr}\!\left(\sigma_{b_R^\alpha D'}^{(\alpha)}\right).
\]
Each $\sigma_{b_R^\alpha D'}^{(\alpha)}$ is normalized, so its trace is~1. The same computation for $\rho_{ABD}$ yields the same value because the right factor is normalized there as well.
\end{proof}

An approximate version follows from recoverability. If the conditional mutual information satisfies
\[
I(A:D\mid B)_\rho \le \varepsilon,
\]
then there exists a recovery map $\mathcal{R}_{B\to BD}$ such that
\[
\left\|
\rho_{ABD}
-
(\mathrm{id}_A\otimes\mathcal{R}_{B\to BD})(\rho_{AB})
\right\|_1
\le
2\sqrt{\ln 2\,\varepsilon}.
\]
This gives a bounded-error observer-modularity theorem: even when the Markov condition is only approximate, the interior observer state is recoverable with controlled error. The collar width controls $\varepsilon$, and thicker collars give better isolation.

This structure is the algebraic version of the patch-net formalism. The sector labels $\alpha$ are the classical center data (the overlap interface), the blockwise factorization is the algebraic version of the interface projection, and the splice theorem is the algebraic version of gauge-invariant normal-form uniqueness: internal observer data is invariant under compatible environment substitution.


\begin{thebibliography}{9}

\bibitem{OPH}
B.~M\"uller,
\emph{Observers Are All You Need: Observer-Patch Holography as a Fundamental Theory},
2026.
Available at \url{https://oph-book.floatingpragma.io/papers/observers_are_all_you_need.pdf}.

\end{thebibliography}

\end{document}
